\clearpage\thispagestyle{empty}
\pdfbookmark[0]{Quellen zum Brief über Pallas und Juno}{pallasjuno18060708-sources}
\begin{center}\textbf{Quellen zum Brief über Pallas und Juno}\end{center}
\noindent Diese KI-Übertragung in heutiges Deutsch folgt dem Brief Carl Friedrich Gauß’ vom 8. Juli 1806 an Franz Xaver von Zach in \href{https://gdz.sub.uni-goettingen.de/id/PPN236007467}{Gauß’ \textit{Werken}, Band~VI} (Göttingen, 1874), S.~277--278. Er erschien in der \textit{Monatlichen Correspondenz}, Band~XIV, August 1806, Beitrag~XVI, Teil~1, \href{https://www.digitale-sammlungen.de/en/view/bsb10538606?page=201}{S.~187}--\href{https://www.digitale-sammlungen.de/en/view/bsb10538606?page=206}{192}. Die Zeitschriftendigitalisate stammen von der Bayerischen Staatsbibliothek.
\par\medskip
Die Werkausgabe gibt in eckigen Klammern eine Beschreibung der Pallas-Ephemeride, die in der Zeitschrift auf \href{https://www.digitale-sammlungen.de/en/view/bsb10538606?page=203}{S.~189}--\href{https://www.digitale-sammlungen.de/en/view/bsb10538606?page=206}{192} abgedruckt ist. Die zitierten Mailänder Beobachtungen stehen auf \href{https://www.digitale-sammlungen.de/en/view/bsb10538606?page=100}{S.~90 des Julihefts}. Dort wird Carlini für die Beobachtungen und ihre Berechnung genannt. Friedrich Wilhelm Bessels Brief und seine Beobachtungen der Juno beginnen auf \href{https://www.digitale-sammlungen.de/en/view/bsb10538606?page=206}{S.~192 des Augusthefts}. Übertragung und Satz wurden mit KI erstellt.
\clearpage
